\section{Implementation details}
\label{app:impl}

\Cref{tab:hparams} lists the training settings.
Both models train for one epoch on \NTrain{} rows, which is 406 steps at a batch of 32.

\begin{table}[h]
  \centering\small
  \caption{Training settings. Learning rates are in \cref{sec:protocol}.}
  \label{tab:hparams}
  \begin{tabular}{@{}lll@{}}
    \toprule
    & UltraBERT-380M & Qwen3-0.6B (full) \\
    \midrule
    Optimiser & AdamW, $\beta=(0.9, 0.98)$, $\epsilon=10^{-6}$ & AdamW \\
    Weight decay & 0.1 & 0.01 \\
    Schedule & 5\% warmup, linear decay to 0 & 100 warmup steps, linear decay to 10\% \\
    Gradient clipping & 1.0 & 1.0 \\
    Precision & fp32 weights, bf16 autocast & fp32 weights, bf16 autocast \\
    Attention & FlashAttention-2, varlen & default \\
    Gradient checkpointing & yes & -- \\
    Token limit & 32{,}768 & 16{,}384 \\
    \bottomrule
  \end{tabular}
\end{table}

The encoder settings are those of the UltraBERT model card, and we change only the learning rate.
No training row exceeds either token limit, because the data are cut at \MaxTokens{} tokens (\cref{sec:data}).

\section{Evidence head}
\label{app:evidence}

The evidence head marks the document tokens that support an answer, and runs on top of a frozen encoder.
We train it on a frozen UltraBERT-380M model, so the decisions of the encoder do not change and one forward pass gives both the decision and the evidence.
We train one head on the model fine-tuned with the sigmoid-normalised loss and one on the model fine-tuned with the softmax, both at a learning rate of $1.5\cdot10^{-5}$, because AT2 and ContextCite were run on the softmax model.

\paragraph{Model.}
For document token $t$ and option $k$ the head reads $[h_t;\,h_k;\,h_t\odot h_k]$, where $h_t$ is the final hidden state of the token and $h_k$ that of option $k$'s mask.
Both come from the one forward pass of the encoder, which reads the question and all the options together with the document, so every $h_t$ already depends on them.
A two-layer perceptron with 512 hidden units, GELU and dropout 0.1 maps this vector to one logit, and its sigmoid is the probability that the token is evidence for option $k$.
The score of a sentence is the highest probability among its tokens.

\paragraph{Labels.}
A token is positive if it overlaps an evidence span.
Where a dataset has gold evidence we use it (Nimble, ContractNLI, MuSiQue, ProofWriter).
For rows with no gold evidence we use generated spans, produced as in \cref{fig:spans}.
Qwen3.6-27B, with thinking, is prompted with the numbered sentences of the row, the question and the gold answer, and asked for the smallest set of sentences that supports the answer and without which it is no longer supported.
We sample ten replies at temperature 1.2 and top-$p$ 0.9, merge identical sets, and drop sets that cite more than 384 tokens unless they are a single sentence.
Each remaining set $S$ is scored by the same model on the decision prompt: once with the full context, once without $S$ and once with only $S$, which makes $1+2m$ prompts for $m$ sets.
The score of each option is the sum of the log-probabilities of its tokens, as in the trie of \cref{sec:loss}, and $q(\mathrm{gold}\mid\cdot)$ is the softmax over the options.
The reward of a set is $\log q(\mathrm{gold}\mid\text{only }S)-\log q(\mathrm{gold}\mid\text{without }S)$, the sum of its sufficiency and its necessity with equal weight and no penalty for size, since the token limit already bounds it.
The set with the highest reward is the label of the row.
The labels cover ANLI, QuALITY, LogiQA~2.0, CosmosQA and the rows of ContractNLI, MuSiQue and ProofWriter without gold evidence.
Unlike SelfCite, which scores the model's own response, we score the gold answer, so the reward carries information only where the labelling model believes the gold answer.
We did not check the generated labels against human evidence on their own.
The indirect check is the ablation below: adding them to the gold labels raises the agreement of the head with human evidence.
Gold evidence is used when both exist.
There are 7{,}778 labelled training rows, of which 7{,}390 are used for training and 388 for validation, and 5.2\% of the training tokens are positive.

\paragraph{Training.}
The loss is token-level binary cross-entropy, with the positive class weighted by the ratio of negative to positive tokens (18.1).
In training the head is conditioned on the gold option, and at test time on the option the encoder chooses.
We use AdamW with a learning rate of $10^{-3}$ and weight decay 0.01, 16 rows per step and three epochs, and keep the epoch with the best token average precision on the validation rows.
The encoder stays frozen.

\paragraph{Evaluation.}
We use the 50 rows per dataset, from ContractNLI, MuSiQue, Nimble and ProofWriter, that have gold evidence and at least two sentences, the same rows on which AT2 and ContextCite are scored.
A word is highlighted when its probability is at least a global threshold, the one that gave the best word F1 on 956 other held-out rows: \EvThrSig{} for the head on the sigmoid-normalised model and \EvThrList{} for the head on the softmax model.
Adding the generated labels to the gold ones raises the sentence average precision from \EvSentAPGold{} to \EvSentAPList{} and the token average precision from \EvTokAPGold{} to \EvTokAPList{} (head on the softmax model).

\section{Additional results}
\label{app:extra}

\Cref{tab:perdataset} gives the test accuracy for each dataset, averaged over seeds.

\begin{table}[h]
  \centering\small
  \setlength{\tabcolsep}{4pt}
  \caption{Test accuracy (\%) per dataset, averaged over seeds. UB is UltraBERT-380M. Bold: best below 1B parameters.}
  \label{tab:perdataset}
  % Generated by paper/figures/make_figures.py; do not edit.
\begin{tabular}{@{}lrccccc@{}}
\toprule
Dataset & Rows & UB sig-norm & UB softmax & Qwen full & Qwen LoRA & Qwen-4B LoRA \\
\midrule
CLINC150 & 400 & \textbf{91.8} & \textbf{91.8} & 86.3 & 79.8 & 89.8 \\
MASSIVE & 400 & \textbf{84.9} & 83.7 & 76.9 & 68.8 & 83.5 \\
Bitext & 400 & \textbf{99.3} & \textbf{99.3} & 96.0 & 94.8 & 97.2 \\
\midrule
ReClor & 400 & \textbf{59.9} & 58.9 & 57.7 & 51.2 & 72.2 \\
LogiQA 2.0 & 400 & 54.8 & 52.9 & \textbf{57.0} & 56.0 & 73.5 \\
CosmosQA & 400 & \textbf{75.2} & 74.8 & 59.9 & 55.0 & 76.0 \\
PIQA & 400 & \textbf{75.1} & 74.8 & 70.2 & 68.8 & 83.8 \\
ANLI & 400 & \textbf{73.6} & 73.1 & 63.7 & 61.0 & 74.0 \\
ProofWriter & 400 & 61.6 & 60.7 & \textbf{63.3} & 63.0 & 70.5 \\
\midrule
ContractNLI & 393 & \textbf{78.6} & 73.8 & 76.1 & 65.6 & 79.1 \\
MuSiQue & 400 & \textbf{69.8} & 58.9 & 62.9 & 54.8 & 71.5 \\
QuALITY & 398 & \textbf{61.3} & 56.9 & 51.5 & 50.0 & 68.8 \\
Nimble & 224 & 50.7 & 47.0 & \textbf{57.3} & 50.9 & 64.3 \\
\bottomrule
\end{tabular}

\end{table}

\section{Examples}
\label{app:examples}

The examples show the decision of the sigmoid-normalised encoder and the token scores of its evidence head for one row of each of three datasets.
The rows are the first of each dataset, in the evaluation order of the evidence study, that has at most 1{,}400 tokens; MuSiQue has none.
Only the tokens that the head scores at or above its threshold of \EvThrSig{} for the chosen option are highlighted, darker for a higher probability, and tokens inside the gold evidence are underlined.
The encoder chooses the correct option in all three rows; the rows were picked by the rule above and not by outcome.

% Generated by paper/figures/make_examples.py; do not edit.
\paragraph{ContractNLI.} 1{,}376 tokens.
The model chooses option 1, which is \textbf{correct}. Question: Statement: "Receiving Party may share some Confidential Information with some third-parties (including consultants, agents and professional advisors)." Does the contract entail the statement, contradict it, or not mention it?

\begin{tabular}{@{}rrp{0.7\linewidth}l@{}}
\toprule
$k$ & Probability & Option & \\
\midrule
1 & 49\% & The contract entails the statement. & gold, chosen \\
2 & 32\% & The contract contradicts the statement. &  \\
3 & 19\% & The contract does not mention it. &  \\
\bottomrule
\end{tabular}

\begingroup\scriptsize\setlength{\parindent}{0pt}\setlength{\parskip}{0.3em}\raggedright
\textsf{\textbf{Contract}}\quad  CONFIDENTIALITY \& NON-DISCLOSURE AGREEMENT ATTACHMENT 3 - INVITATION FOR BID RFQ 205 CONSTRUCTION OF THE ASHBURN SHERIFF STATION THIS AGREEMENT is entered into by and between COUNTY OF LOUDOUN, VIRGINIA with offices at 1 Harrison St. S.E., Leesburg, Virginia 20175 (hereinafter ``County'') and with offices at (hereinafter ``Bidder'') on \_\_\_\_\_\_\_\_\_\_\_\_\_\_\_\_\_\_\_\_\_2015. Whereas, the County is seeking bids from Bidders to construct the Ashburn Sheriff Station (the ``Project''); and Whereas, the Bidder desires to obtain architectural {\fboxsep=0pt\colorbox[HTML]{A8C8EE}{drawings}} that detail the design of said project for the sole purpose of compiling and submitting a response to Invitation for Bid RFQ 205 to construct the Project. NOW IT IS THEREFORE AGREED: 1. Purpose The purpose of this agreement is to allow the Bidder access to Confidential Information as defined herein detailing {\fboxsep=0pt\colorbox[HTML]{BFD7F3}{the}} design of the Project for the sole purpose of compiling and submitting a bid in response to the County's Invitation for Bid to construct the Project. This agreement is intended to ensure the safety and security of the Ashburn Sheriff Station. By signing this agreement the Bidder is stating that it will keep all Confidential Information (as defined in paragraph 4) {\fboxsep=0pt\colorbox[HTML]{C1D8F3}{in}} {\fboxsep=0pt\colorbox[HTML]{B4CFF0}{strict}} confidence and only use the Confidential Information on the Project identified herein. Further the Bidder warrants that its {\fboxsep=0pt\colorbox[HTML]{BED6F2}{employees}}, partners, subcontractors, agents, and clients (collectively hereinafter the ``Bidder'') will be similarly bound by this agreement, or by other similar constraints of confidentiality. 2. Disclosure \underline{The} \underline{B}\underline{id}\underline{der} \underline{shall} {\fboxsep=0pt\colorbox[HTML]{BBD4F2}{\underline{not}}} \underline{publish}\underline{,} \underline{or} \underline{make} \underline{and}\underline{/}\underline{or} {\fboxsep=0pt\colorbox[HTML]{C0D7F3}{\underline{retain}}} \underline{copies} \underline{or} \underline{disclose} \underline{to} \underline{any} \underline{third} \underline{party}\underline{,} \underline{any} \underline{Conf}\underline{idential} \underline{Information} \underline{furnished} \underline{to} \underline{the} \underline{B}\underline{id}\underline{der} \underline{except} \underline{to} \underline{such} \underline{extent} \underline{as} \underline{may} \underline{be} \underline{necessary} \underline{to} \underline{carry} \underline{out} \underline{the} \underline{B}\underline{id}\underline{der}\underline{'}\underline{s} \underline{duties} \underline{under} \underline{this} \underline{agreement} \underline{and} {\fboxsep=0pt\colorbox[HTML]{B9D2F1}{\underline{except}}} \underline{to} \underline{the} \underline{extent} \underline{that} \underline{disclosure} \underline{is} \underline{required} \underline{by} \underline{judicial} \underline{process} \underline{in} \underline{a} \underline{legal} \underline{proceeding}\underline{,} \underline{in} \underline{which} \underline{event} \underline{B}\underline{id}\underline{der} \underline{will} \underline{provide} \underline{to} \underline{the} \underline{County} \underline{notice} \underline{of} \underline{the} \underline{request} \underline{or} \underline{demand} \underline{to} \underline{disclose} \underline{Conf}\underline{idential} \underline{Information}\underline{.} Within ten (10) days after award and/or completion of the IFB, or upon demand by the County, all Confidential Information, including all photographs, sketches, models, and drawings related to the design, layout and operation of said project shall be destroyed and files deleted from all of the Bidder's web sites and electronic devices (including but not limited to computers, cellular phones, etc.). Bidder shall also verify that all Confidential Information provided to partners, subcontractors, agents and clients has been destroyed. The successful Bidder shall destroy all Confidential Information within ten (10) days after final Project completion. 3. Indemnification Bidder agrees to indemnify and hold harmless the County, its elected officials, employees and agents from and against any and all liability, damages, claims, suits, liens, and judgments (including reasonable attorney's fees), of whatever nature, for injuries to or death of any person or persons, or loss of or damage to property, to the extent attributable to the negligent acts of Bidder, its subcontractors or its respective agents, servants, or employees or such parties' related to said Confidential Information. 4. Confidentiality A. ``Confidential Information'' as outlined herein shall be used in {\fboxsep=0pt\colorbox[HTML]{AFCCF0}{strict}} confidence solely for the purpose allotted under this agreement and shall {\fboxsep=0pt\colorbox[HTML]{BBD4F2}{not}} {\fboxsep=0pt\colorbox[HTML]{9BC0EC}{be}} disclosed by the Bidder{\fboxsep=0pt\colorbox[HTML]{A2C4ED}{,}} its {\fboxsep=0pt\colorbox[HTML]{C1D8F3}{employees}}, partners, subcontractors, agents, or clients to a third party without the prior written consent of the County with the exception of the following; \textbullet{} It was already in the possession of the receiving party; \textbullet{} It is learned from or becomes publicly disclosed by a source with legal right to use and disclose same \textbullet{} It is identical to information which was developed by the receiving party independently of any Confidential Information received from the other party. B. Confidential Information means all information relating to the design, layout and operation of said Project. 5. Governing Law This agreement validity, construction and effect shall be governed under the substantive law of the Commonwealth of Virginia without the giving effect to any choice of law provision. Bidder expressly waives any objection to venue or jurisdiction of the Loudoun County Circuit Court, Loudoun County, Virginia. 6. Entire Agreement The governing terms and conditions of this agreement are expressly limited to the terms and conditions contained in this Agreement. This agreement constitutes the complete integration of all oral and written documents, is the entire and final agreement between the parties and may be amended only in writing signed by authorized officials of both parties. 7. Authority to Bind Firm in Agreement Person signing Agreement will show title or authority to bind the firm in an agreement. Name and authorized signature must appear below. Those authorized to sign are as follows: If a sole proprietorship, the owner may sign. If a general partnership, any general partner may sign. If a limited partnership, a general partner must sign. If a limited liability company, a ``member'' may sign or ``manager'' must sign if so specified by the articles or organization. If a regular corporation, the CEO, President or Vice-President must sign. Others may be granted authority to sign but the County requires that a corporate document authorizing him/her to sign be submitted with proposal. AGREED AND ACCEPTED BY: By: \_\_\_\_\_\_\_\_\_\_\_\_\_\_\_\_\_\_\_\_\_\_\_\_\_\_\_\_\_\_\_\_\_\_\_\_\_\_\_\_\_ Name (printed): \_\_\_\_\_\_\_\_\_\_\_\_\_\_\_\_\_\_\_\_\_\_\_\_\_\_\_\_\_\_\_ Title: \_\_\_\_\_\_\_\_\_\_\_\_\_\_\_\_\_\_\_\_\_\_\_\_\_\_\_\_\_\_\_\_\_\_\_\_\_\_\_\_ Date: \_\_\_\_\_\_\_\_\_\_\_\_\_\_\_\_\_\_\_\_\_\_\_\_\_\_\_\_\_\_\_\_\_\_\_\_\_\_\_ COMMONWEALTH OF VIRGINIA COUNTY OF \_\_\_\_\_\_\_\_\_\_\_\_\_\_\_\_\_\_\_\_ to-wit: Acknowledged before me this \_\_\_\_ day of \_\_\_\_\_\_\_\_\_\_\_\_\_\_\_\_\_\_\_\_\_\_\_, 2015. Notary Public Registration No.: My Commission Expires: \_\_\_\_\_\_\_\_\_\_\_\_\_\_\_\_\_\_\_\_\_\_\_\_ 

\endgroup

\paragraph{Nimble.} 479 tokens.
The model chooses option 2, which is \textbf{correct}. Question: Is the supplier bid ready to route for approval now under the stated rubric? Resolve pronouns using their nearest coherent antecedent.

\begin{tabular}{@{}rrp{0.7\linewidth}l@{}}
\toprule
$k$ & Probability & Option & \\
\midrule
1 & 46\% & No --- at least one stated requirement is unsatisfied now, so do not route for approval; route for supplier clarification. &  \\
2 & 54\% & Yes --- every commercial, capacity, certification, technical, and current attachment requirement is satisfied, so route for approval. & gold, chosen \\
\bottomrule
\end{tabular}

\begingroup\scriptsize\setlength{\parindent}{0pt}\setlength{\parskip}{0.3em}\raggedright
\textsf{\textbf{Statements}}\quad Hand{\fboxsep=0pt\colorbox[HTML]{BED6F3}{off}} coordinator: {\fboxsep=0pt\colorbox[HTML]{7AABE5}{The}} {\fboxsep=0pt\colorbox[HTML]{8CB6E9}{PR}}{\fboxsep=0pt\colorbox[HTML]{77A9E5}{-}}{\fboxsep=0pt\colorbox[HTML]{7CACE6}{4}}{\fboxsep=0pt\colorbox[HTML]{6DA3E3}{8}}{\fboxsep=0pt\colorbox[HTML]{74A7E4}{2}} package {\fboxsep=0pt\colorbox[HTML]{99BEEB}{was}} {\fboxsep=0pt\colorbox[HTML]{A1C4ED}{evaluated}} {\fboxsep=0pt\colorbox[HTML]{96BDEB}{at}} {\fboxsep=0pt\colorbox[HTML]{669EE2}{the}} {\fboxsep=0pt\colorbox[HTML]{9CC0EC}{2}}{\fboxsep=0pt\colorbox[HTML]{659DE1}{0}}{\fboxsep=0pt\colorbox[HTML]{689FE2}{2}}{\fboxsep=0pt\colorbox[HTML]{639CE1}{6}}{\fboxsep=0pt\colorbox[HTML]{679FE2}{-}}{\fboxsep=0pt\colorbox[HTML]{9ABFEB}{0}}{\fboxsep=0pt\colorbox[HTML]{659EE1}{9}}{\fboxsep=0pt\colorbox[HTML]{679FE2}{-}}{\fboxsep=0pt\colorbox[HTML]{669EE2}{1}}{\fboxsep=0pt\colorbox[HTML]{69A0E2}{7}} {\fboxsep=0pt\colorbox[HTML]{89B4E8}{operational}} {\fboxsep=0pt\colorbox[HTML]{81AFE7}{hand}}{\fboxsep=0pt\colorbox[HTML]{8DB7E9}{off}} using the {\fboxsep=0pt\colorbox[HTML]{81AFE7}{quotation}} and documents then attached{\fboxsep=0pt\colorbox[HTML]{9BBFEC}{.}} Commercial reviewer: {\fboxsep=0pt\colorbox[HTML]{9ABFEC}{The}} {\fboxsep=0pt\colorbox[HTML]{B0CDF0}{supplier}} {\fboxsep=0pt\colorbox[HTML]{639CE1}{quoted}} {\fboxsep=0pt\colorbox[HTML]{92BAEA}{a}} {\fboxsep=0pt\colorbox[HTML]{ADCBEF}{total}} {\fboxsep=0pt\colorbox[HTML]{B7D1F1}{price}} {\fboxsep=0pt\colorbox[HTML]{73A6E4}{of}} {\fboxsep=0pt\colorbox[HTML]{C1D8F3}{\$}}{\fboxsep=0pt\colorbox[HTML]{A1C3ED}{1}}{\fboxsep=0pt\colorbox[HTML]{7BABE6}{1}}{\fboxsep=0pt\colorbox[HTML]{C4DAF4}{,}}{\fboxsep=0pt\colorbox[HTML]{71A5E4}{8}}{\fboxsep=0pt\colorbox[HTML]{71A5E4}{5}}{\fboxsep=0pt\colorbox[HTML]{9BC0EC}{0}} and {\fboxsep=0pt\colorbox[HTML]{8CB6E9}{a}} {\fboxsep=0pt\colorbox[HTML]{87B3E8}{lead}} {\fboxsep=0pt\colorbox[HTML]{AFCCF0}{time}} {\fboxsep=0pt\colorbox[HTML]{77A9E5}{of}} 2{\fboxsep=0pt\colorbox[HTML]{BED6F2}{8}} {\fboxsep=0pt\colorbox[HTML]{82B0E7}{calendar}} days {\fboxsep=0pt\colorbox[HTML]{649DE1}{for}} {\fboxsep=0pt\colorbox[HTML]{639CE1}{PR}}{\fboxsep=0pt\colorbox[HTML]{6DA3E3}{-}}{\fboxsep=0pt\colorbox[HTML]{6CA2E3}{4}}{\fboxsep=0pt\colorbox[HTML]{6BA1E3}{8}}{\fboxsep=0pt\colorbox[HTML]{689FE2}{2}}{\fboxsep=0pt\colorbox[HTML]{8BB5E9}{.}} Quotation reviewer: {\fboxsep=0pt\colorbox[HTML]{7CACE6}{\underline{At}}} {\fboxsep=0pt\colorbox[HTML]{649DE1}{\underline{the}}} {\fboxsep=0pt\colorbox[HTML]{A4C6EE}{\underline{2}}}{\fboxsep=0pt\colorbox[HTML]{639CE1}{\underline{0}}}{\fboxsep=0pt\colorbox[HTML]{629CE1}{\underline{2}}}{\fboxsep=0pt\colorbox[HTML]{669EE2}{\underline{6}}}{\fboxsep=0pt\colorbox[HTML]{679FE2}{\underline{-}}}{\fboxsep=0pt\colorbox[HTML]{9EC2EC}{\underline{0}}}{\fboxsep=0pt\colorbox[HTML]{639CE1}{\underline{9}}}{\fboxsep=0pt\colorbox[HTML]{95BCEB}{\underline{-}}}{\fboxsep=0pt\colorbox[HTML]{6FA4E3}{\underline{1}}}{\fboxsep=0pt\colorbox[HTML]{669EE2}{\underline{7}}} {\fboxsep=0pt\colorbox[HTML]{6BA1E3}{\underline{operational}}} {\fboxsep=0pt\colorbox[HTML]{659DE1}{\underline{hand}}}{\fboxsep=0pt\colorbox[HTML]{619BE1}{\underline{off}}}{\fboxsep=0pt\colorbox[HTML]{679EE2}{\underline{,}}} {\fboxsep=0pt\colorbox[HTML]{669EE2}{\underline{the}}} {\fboxsep=0pt\colorbox[HTML]{83B1E7}{\underline{supplier}}} {\fboxsep=0pt\colorbox[HTML]{6DA3E3}{\underline{quotation}}} {\fboxsep=0pt\colorbox[HTML]{669EE2}{\underline{for}}} {\fboxsep=0pt\colorbox[HTML]{619BE1}{\underline{PR}}}{\fboxsep=0pt\colorbox[HTML]{6DA2E3}{\underline{-}}}{\fboxsep=0pt\colorbox[HTML]{6CA2E3}{\underline{4}}}{\fboxsep=0pt\colorbox[HTML]{73A6E4}{\underline{8}}}{\fboxsep=0pt\colorbox[HTML]{669EE2}{\underline{2}}} {\fboxsep=0pt\colorbox[HTML]{87B3E8}{\underline{stated}}} {\fboxsep=0pt\colorbox[HTML]{7FAEE6}{\underline{a}}} {\fboxsep=0pt\colorbox[HTML]{80AFE7}{\underline{monthly}}} {\fboxsep=0pt\colorbox[HTML]{C7DBF4}{\underline{g}}}{\fboxsep=0pt\colorbox[HTML]{B0CDF0}{\underline{asket}}} {\fboxsep=0pt\colorbox[HTML]{72A6E4}{\underline{capacity}}} {\fboxsep=0pt\colorbox[HTML]{659DE1}{\underline{of}}} {\fboxsep=0pt\colorbox[HTML]{6FA4E3}{\underline{6}}}{\fboxsep=0pt\colorbox[HTML]{70A4E3}{\underline{4}}}{\fboxsep=0pt\colorbox[HTML]{74A7E4}{\underline{0}}} {\fboxsep=0pt\colorbox[HTML]{9EC2EC}{\underline{g}}}{\fboxsep=0pt\colorbox[HTML]{8FB8E9}{\underline{ask}}}{\fboxsep=0pt\colorbox[HTML]{A4C5ED}{\underline{ets}}}{\fboxsep=0pt\colorbox[HTML]{84B1E7}{\underline{.}}} Request reviewer: {\fboxsep=0pt\colorbox[HTML]{78A9E5}{\underline{At}}} {\fboxsep=0pt\colorbox[HTML]{6BA1E3}{\underline{the}}} {\fboxsep=0pt\colorbox[HTML]{A8C8EE}{\underline{2}}}{\fboxsep=0pt\colorbox[HTML]{639CE1}{\underline{0}}}{\fboxsep=0pt\colorbox[HTML]{639CE1}{\underline{2}}}{\fboxsep=0pt\colorbox[HTML]{669EE1}{\underline{6}}}{\fboxsep=0pt\colorbox[HTML]{639CE1}{\underline{-}}}{\fboxsep=0pt\colorbox[HTML]{9BC0EC}{\underline{0}}}{\fboxsep=0pt\colorbox[HTML]{639CE1}{\underline{9}}}{\fboxsep=0pt\colorbox[HTML]{8AB5E8}{\underline{-}}}{\fboxsep=0pt\colorbox[HTML]{B5D0F1}{\underline{1}}}{\fboxsep=0pt\colorbox[HTML]{669EE2}{\underline{7}}} {\fboxsep=0pt\colorbox[HTML]{6EA3E3}{\underline{operational}}} {\fboxsep=0pt\colorbox[HTML]{8AB5E8}{\underline{hand}}}{\fboxsep=0pt\colorbox[HTML]{609AE0}{\underline{off}}}{\fboxsep=0pt\colorbox[HTML]{72A6E4}{\underline{,}}} {\fboxsep=0pt\colorbox[HTML]{71A5E4}{\underline{PR}}}{\fboxsep=0pt\colorbox[HTML]{7AAAE5}{\underline{-}}}{\fboxsep=0pt\colorbox[HTML]{72A6E4}{\underline{4}}}{\fboxsep=0pt\colorbox[HTML]{6DA3E3}{\underline{8}}}{\fboxsep=0pt\colorbox[HTML]{6EA3E3}{\underline{2}}} {\fboxsep=0pt\colorbox[HTML]{A8C8EE}{\underline{stated}}} \underline{a} {\fboxsep=0pt\colorbox[HTML]{A8C8EE}{\underline{minimum}}} {\fboxsep=0pt\colorbox[HTML]{78A9E5}{\underline{monthly}}} {\fboxsep=0pt\colorbox[HTML]{B0CDF0}{\underline{g}}}{\fboxsep=0pt\colorbox[HTML]{A6C6EE}{\underline{asket}}} {\fboxsep=0pt\colorbox[HTML]{6CA2E3}{\underline{capacity}}} {\fboxsep=0pt\colorbox[HTML]{8EB7E9}{\underline{of}}} {\fboxsep=0pt\colorbox[HTML]{7AABE5}{\underline{6}}}{\fboxsep=0pt\colorbox[HTML]{8BB5E9}{\underline{0}}}{\fboxsep=0pt\colorbox[HTML]{93BAEA}{\underline{0}}} {\fboxsep=0pt\colorbox[HTML]{81AFE7}{\underline{g}}}{\fboxsep=0pt\colorbox[HTML]{89B4E8}{\underline{ask}}}{\fboxsep=0pt\colorbox[HTML]{ACCAEF}{\underline{ets}}}{\fboxsep=0pt\colorbox[HTML]{87B3E8}{\underline{.}}} Certification reviewer: {\fboxsep=0pt\colorbox[HTML]{99BEEB}{The}} {\fboxsep=0pt\colorbox[HTML]{89B4E8}{supplier}}{\fboxsep=0pt\colorbox[HTML]{8CB6E9}{'s}} {\fboxsep=0pt\colorbox[HTML]{C2D9F3}{ISO}} 9{\fboxsep=0pt\colorbox[HTML]{BBD4F2}{0}}01 {\fboxsep=0pt\colorbox[HTML]{BFD7F3}{certification}} was current when {\fboxsep=0pt\colorbox[HTML]{9DC1EC}{the}} package was {\fboxsep=0pt\colorbox[HTML]{8EB7E9}{evaluated}}{\fboxsep=0pt\colorbox[HTML]{A6C7EE}{,}} and its {\fboxsep=0pt\colorbox[HTML]{95BCEB}{current}} {\fboxsep=0pt\colorbox[HTML]{AFCCF0}{certificate}} was {\fboxsep=0pt\colorbox[HTML]{86B2E8}{attached}} {\fboxsep=0pt\colorbox[HTML]{71A5E4}{to}} {\fboxsep=0pt\colorbox[HTML]{6DA2E3}{the}} {\fboxsep=0pt\colorbox[HTML]{629CE1}{PR}}{\fboxsep=0pt\colorbox[HTML]{A2C4ED}{-}}{\fboxsep=0pt\colorbox[HTML]{77A9E5}{4}}{\fboxsep=0pt\colorbox[HTML]{88B4E8}{8}}{\fboxsep=0pt\colorbox[HTML]{669EE2}{2}} {\fboxsep=0pt\colorbox[HTML]{8CB6E9}{package}} at {\fboxsep=0pt\colorbox[HTML]{C4DAF4}{that}} {\fboxsep=0pt\colorbox[HTML]{649DE1}{time}}{\fboxsep=0pt\colorbox[HTML]{A2C4ED}{.}} Document reviewer: The required material {\fboxsep=0pt\colorbox[HTML]{8FB8E9}{certificate}} {\fboxsep=0pt\colorbox[HTML]{99BEEB}{was}} {\fboxsep=0pt\colorbox[HTML]{7FAEE6}{attached}} {\fboxsep=0pt\colorbox[HTML]{6CA2E3}{to}} {\fboxsep=0pt\colorbox[HTML]{679FE2}{the}} {\fboxsep=0pt\colorbox[HTML]{629CE1}{PR}}{\fboxsep=0pt\colorbox[HTML]{8AB5E8}{-}}{\fboxsep=0pt\colorbox[HTML]{83B0E7}{4}}{\fboxsep=0pt\colorbox[HTML]{9EC2EC}{8}}{\fboxsep=0pt\colorbox[HTML]{8BB6E9}{2}} {\fboxsep=0pt\colorbox[HTML]{90B9EA}{package}} when {\fboxsep=0pt\colorbox[HTML]{91B9EA}{evaluated}}. The {\fboxsep=0pt\colorbox[HTML]{98BEEB}{submitted}} {\fboxsep=0pt\colorbox[HTML]{BCD5F2}{drawing}} {\fboxsep=0pt\colorbox[HTML]{8DB7E9}{bore}} {\fboxsep=0pt\colorbox[HTML]{97BDEB}{the}} required signature and was also {\fboxsep=0pt\colorbox[HTML]{C1D8F3}{attached}} then{\fboxsep=0pt\colorbox[HTML]{B8D2F1}{.}} Technical reviewer: The g{\fboxsep=0pt\colorbox[HTML]{AFCCF0}{asket}} {\fboxsep=0pt\colorbox[HTML]{679FE2}{specification}} {\fboxsep=0pt\colorbox[HTML]{ACCAEF}{in}} {\fboxsep=0pt\colorbox[HTML]{80AEE7}{the}} {\fboxsep=0pt\colorbox[HTML]{A6C7EE}{supplier}} {\fboxsep=0pt\colorbox[HTML]{689FE2}{quotation}} matches the technical {\fboxsep=0pt\colorbox[HTML]{75A7E4}{specification}} {\fboxsep=0pt\colorbox[HTML]{98BEEB}{in}} {\fboxsep=0pt\colorbox[HTML]{71A5E4}{PR}}{\fboxsep=0pt\colorbox[HTML]{77A9E5}{-}}{\fboxsep=0pt\colorbox[HTML]{6AA0E2}{4}}{\fboxsep=0pt\colorbox[HTML]{82B0E7}{8}}{\fboxsep=0pt\colorbox[HTML]{7EAEE6}{2}}{\fboxsep=0pt\colorbox[HTML]{7BABE6}{.}} Pro{\fboxsep=0pt\colorbox[HTML]{81AFE7}{c}}urement specialist: Approval requires price no more than \$12,0{\fboxsep=0pt\colorbox[HTML]{9FC2ED}{0}}{\fboxsep=0pt\colorbox[HTML]{99BEEB}{0}}, {\fboxsep=0pt\colorbox[HTML]{8CB6E9}{lead}} time no more than 30 {\fboxsep=0pt\colorbox[HTML]{95BCEB}{calendar}} days, {\fboxsep=0pt\colorbox[HTML]{98BEEB}{capacity}} at least {\fboxsep=0pt\colorbox[HTML]{7EAEE6}{6}}{\fboxsep=0pt\colorbox[HTML]{A5C6EE}{0}}{\fboxsep=0pt\colorbox[HTML]{C3D9F3}{0}}/{\fboxsep=0pt\colorbox[HTML]{80AFE7}{month}}, current ISO 9{\fboxsep=0pt\colorbox[HTML]{8EB7E9}{0}}{\fboxsep=0pt\colorbox[HTML]{8FB8E9}{0}}1, a {\fboxsep=0pt\colorbox[HTML]{A8C8EE}{material}} certificate{\fboxsep=0pt\colorbox[HTML]{AFCDF0}{,}} and a {\fboxsep=0pt\colorbox[HTML]{B7D1F1}{signed}} {\fboxsep=0pt\colorbox[HTML]{B4CFF1}{drawing}}---all documents attached now. Technical {\fboxsep=0pt\colorbox[HTML]{C1D8F3}{specification}} owner: technical match does not {\fboxsep=0pt\colorbox[HTML]{9ABFEC}{wa}}ive the attachment requirements{\fboxsep=0pt\colorbox[HTML]{90B8EA}{.}} P{\fboxsep=0pt\colorbox[HTML]{A1C4ED}{urch}}asing approver: Packages meeting every threshold{\fboxsep=0pt\colorbox[HTML]{C1D7F3}{,}} including boundary values, may {\fboxsep=0pt\colorbox[HTML]{7FAEE6}{be}} approved. If any required document {\fboxsep=0pt\colorbox[HTML]{86B2E8}{is}} not attached now, route the package for supplier clarification{\fboxsep=0pt\colorbox[HTML]{A4C5EE}{.}}

\endgroup

\paragraph{ProofWriter.} 195 tokens.
The model chooses option 2, which is \textbf{correct}. Question: Statement: "Gary is not blue." Based only on the theory, is the statement true, false, or unknown?

\begin{tabular}{@{}rrp{0.7\linewidth}l@{}}
\toprule
$k$ & Probability & Option & \\
\midrule
1 & 1\% & True: the theory implies the statement. &  \\
2 & 69\% & False: the theory implies that the statement is false. & gold, chosen \\
3 & 30\% & Unknown: the theory implies neither the statement nor its negation. &  \\
\bottomrule
\end{tabular}

\begingroup\scriptsize\setlength{\parindent}{0pt}\setlength{\parskip}{0.3em}\raggedright
\textsf{\textbf{Theory}}\quad Bob is {\fboxsep=0pt\colorbox[HTML]{C6DBF4}{big}}{\fboxsep=0pt\colorbox[HTML]{B7D1F1}{.}} Bob is blue. Bob is cold. Bob is white. Dave is cold. Dave is kind. Dave {\fboxsep=0pt\colorbox[HTML]{C7DCF4}{is}} nice{\fboxsep=0pt\colorbox[HTML]{C0D7F3}{.}} {\fboxsep=0pt\colorbox[HTML]{7FAEE6}{Gary}} {\fboxsep=0pt\colorbox[HTML]{88B3E8}{is}} {\fboxsep=0pt\colorbox[HTML]{A0C3ED}{cold}}{\fboxsep=0pt\colorbox[HTML]{C5DAF4}{.}} {\fboxsep=0pt\colorbox[HTML]{C6DBF4}{\underline{Gary}}} \underline{is} {\fboxsep=0pt\colorbox[HTML]{B5D0F1}{\underline{kind}}}\underline{.} Harry is {\fboxsep=0pt\colorbox[HTML]{B7D1F1}{white}}{\fboxsep=0pt\colorbox[HTML]{ACCAEF}{.}} {\fboxsep=0pt\colorbox[HTML]{85B1E7}{\underline{All}}} {\fboxsep=0pt\colorbox[HTML]{84B1E7}{\underline{white}}} {\fboxsep=0pt\colorbox[HTML]{80AFE7}{\underline{things}}} {\fboxsep=0pt\colorbox[HTML]{78AAE5}{\underline{are}}} {\fboxsep=0pt\colorbox[HTML]{B2CEF0}{\underline{blue}}}{\fboxsep=0pt\colorbox[HTML]{B1CEF0}{\underline{.}}} {\fboxsep=0pt\colorbox[HTML]{97BDEB}{All}} big{\fboxsep=0pt\colorbox[HTML]{BED5F2}{,}} {\fboxsep=0pt\colorbox[HTML]{AECCEF}{kind}} {\fboxsep=0pt\colorbox[HTML]{93BBEA}{things}} {\fboxsep=0pt\colorbox[HTML]{90B9EA}{are}} {\fboxsep=0pt\colorbox[HTML]{A3C5ED}{nice}}{\fboxsep=0pt\colorbox[HTML]{BED6F2}{.}} \underline{If} {\fboxsep=0pt\colorbox[HTML]{99BEEB}{\underline{something}}} \underline{is} \underline{nice} \underline{then} \underline{it} {\fboxsep=0pt\colorbox[HTML]{A6C7EE}{\underline{is}}} {\fboxsep=0pt\colorbox[HTML]{B5D0F1}{\underline{big}}}\underline{.} {\fboxsep=0pt\colorbox[HTML]{C3D9F3}{\underline{If}}} {\fboxsep=0pt\colorbox[HTML]{A6C6EE}{\underline{something}}} {\fboxsep=0pt\colorbox[HTML]{ADCBEF}{\underline{is}}} {\fboxsep=0pt\colorbox[HTML]{B0CDF0}{\underline{kind}}} \underline{then} {\fboxsep=0pt\colorbox[HTML]{AFCCF0}{\underline{it}}} {\fboxsep=0pt\colorbox[HTML]{A2C4ED}{\underline{is}}} \underline{nice}\underline{.} {\fboxsep=0pt\colorbox[HTML]{BBD4F2}{\underline{If}}} {\fboxsep=0pt\colorbox[HTML]{B1CEF0}{\underline{something}}} \underline{is} {\fboxsep=0pt\colorbox[HTML]{C2D8F3}{\underline{rough}}} {\fboxsep=0pt\colorbox[HTML]{B9D2F1}{\underline{then}}} {\fboxsep=0pt\colorbox[HTML]{90B9EA}{\underline{it}}} {\fboxsep=0pt\colorbox[HTML]{ACCAEF}{\underline{is}}} \underline{white}{\fboxsep=0pt\colorbox[HTML]{ACCBEF}{\underline{.}}} \underline{If} {\fboxsep=0pt\colorbox[HTML]{A8C8EE}{\underline{something}}} {\fboxsep=0pt\colorbox[HTML]{C6DBF4}{\underline{is}}} \underline{big} \underline{and} {\fboxsep=0pt\colorbox[HTML]{BAD3F2}{\underline{nice}}} \underline{then} {\fboxsep=0pt\colorbox[HTML]{89B4E8}{\underline{it}}} {\fboxsep=0pt\colorbox[HTML]{ABCAEF}{\underline{is}}} {\fboxsep=0pt\colorbox[HTML]{8FB8E9}{\underline{rough}}}{\fboxsep=0pt\colorbox[HTML]{B1CDF0}{\underline{.}}} {\fboxsep=0pt\colorbox[HTML]{8CB6E9}{All}} {\fboxsep=0pt\colorbox[HTML]{9FC2ED}{white}}{\fboxsep=0pt\colorbox[HTML]{A8C8EE}{,}} {\fboxsep=0pt\colorbox[HTML]{B0CDF0}{blue}} {\fboxsep=0pt\colorbox[HTML]{9CC0EC}{things}} {\fboxsep=0pt\colorbox[HTML]{7EADE6}{are}} {\fboxsep=0pt\colorbox[HTML]{92BAEA}{cold}}{\fboxsep=0pt\colorbox[HTML]{A9C8EE}{.}}

\endgroup

